\documentclass[12pt]{article}
\title{Edzőterem működését támogató szoftver PureScript és Vue.js alapokon}
\author{Mezei Botond, Szabó Benedek}
\date{2023}
\renewcommand{\contentsname}{Tartalomjegyzék}
\usepackage[utf8]{inputenc}
\usepackage{graphicx}
\graphicspath{ {./images/} }
\usepackage{wrapfig}
\renewcommand{\figurename}{Ábra}
\usepackage{amsmath}
\usepackage{subcaption}
\usepackage{amssymb}
\usepackage[export]{adjustbox}
\usepackage[magyar]{babel}
\usepackage[breaklinks]{hyperref}
\usepackage{microtype}
\usepackage[T1]{fontenc}
\usepackage{lmodern}
\usepackage{multirow}
\usepackage[table]{xcolor}
%\usepackage{natbib}

\setlength{\arrayrulewidth}{0.5mm}
\setlength{\tabcolsep}{7pt}
\renewcommand{\arraystretch}{2.5}

\textwidth=14cm \textheight=20cm
%\hoffset=-1cm
%\voffset=-1cm

\newtheorem{theorem}{Theorem}[section]
\newtheorem{acknowledgement}[theorem]{Acknowledgement}
\newtheorem{algorithm}[theorem]{Algorithm}
\newtheorem{axiom}[theorem]{Axiom}
\newtheorem{case}[theorem]{Case}
\newtheorem{claim}[theorem]{Claim}
\newtheorem{conclusion}[theorem]{Conclusion}
\newtheorem{condition}[theorem]{Condition}
\newtheorem{conjecture}[theorem]{Conjecture}
\newtheorem{corollary}[theorem]{Corollary}
\newtheorem{criterion}[theorem]{Criterion}
\newtheorem{definíció}[theorem]{Definíció}
\newtheorem{discussion}[theorem]{Discussion}
\newtheorem{example}[theorem]{Example}
\newtheorem{exercise}[theorem]{Exercise}
\newtheorem{explanation}[theorem]{Explanation}
\newtheorem{illustration}[theorem]{Illustration}
\newtheorem{lemma}[theorem]{Lemma}
\newtheorem{notation}[theorem]{Notation}
\newtheorem{problem}[theorem]{Problem}
\newtheorem{proposition}[theorem]{Proposition}
\newtheorem{remark}[theorem]{Remark}
\newtheorem{solution}[theorem]{Solution}
\newtheorem{summary}[theorem]{Summary}

\makeatletter
%\renewcommand*{\@biblabel}[1]{–}
\makeatother

\begin{document}
\begin{titlepage}

	\begin{center}
	\mbox{}\\
	\vspace{35mm}
	\huge
	\textsc{Diplomamunka}
	\end{center}

	
	\vspace{55mm}
	\Large
	\hspace*{\fill} 
	\textbf{Mezei Botond, Szabó Benedek}
	
	\begin{center}
	\vspace{55mm}
	Debrecen\\	
	2023
	
	\end{center}
\end{titlepage}
\begin{titlepage}

	\vspace*{-2cm}
	\hspace*{-1.5cm}

	\begin{center}
	\large
	Debreceni Egyetem
	
	Informatikai Kar	

	Számítógéptudományi Tanszék
	
	\vspace{17mm}
	\huge
	\LARGE
    \textbf{Edzőterem működését támogató szoftver PureScript és Vue.js alapokon}
	
	\vspace{15mm}
	\large
	\textsc{Diplomamunka}
	
	\normalsize
	\vspace{17mm}
	\textsc{Készítette:}
	
	\vspace{5mm}
	\Large
	\textbf{Mezei Botond és Szabó Benedek}
	
	\normalsize
	programtervező informatika szakos hallgatók
	
	\vspace{18mm}
	\textsc{Témavezető:}\\
	\vspace{5mm}
	\large
	Dr. Battyányi Péter \\
	\normalsize
	adjunktus\\
	
	\vspace{27mm}
	Debrecen\\	
	2023
	
	\end{center}
\end{titlepage}
\newpage
\tableofcontents
\newpage
\section*{Bevezetés}
\addcontentsline{toc}{section}{Bevezetés}

A huszonegyedik század harmadik évtizedében minden vállalatnak, cégnek vagy szervezetnek elengedhetetlenül szükségessé vált az informatika bevonása a mindennapi működésbe. A cégek versenyképességük megőrzése, a hatékonyság és profitabilitás növelése érdekében már gyakorlatilag egy mikrovállalkozás méretétől kezdve rászorulnak valamiféle speciális céges szoftver, vállalatirányítási rendszer használatára. Egy kicsit nagyobb vállalatnál pedig már egészen elképzelhetetlen lenne a működés e szoftverek nélkül. A költségcsökkentés és automatizálás mellett meg kell említeni az információ növekvő szerepét, egy vállalatirányítási rendszer biztosíthatja a gyors (valós idejű) információ hozzáférést a munkatársaknak, illetve adatokkal látja el a döntéshozókat.

\begin{figure}
\includegraphics[width=\textwidth]{testimg}
\caption{Az ARIS objektumai, és az ARIS ház felépítése. A jobb oldali ábra forrása:  \textit{Fonó, 2005}}
\end{figure}

\newpage
\section{Szakirodalmi áttekintés}

\subsection{Funkcionális programozás}
Ebben a fejezetben először ismertetem a modellezés motivációit, majd áttekintésre kerül néhány lehetséges modellezési eszköz, melyek üzleti folyamatok, workflow-k modellezése esetén szóba jöhetnek. \cite{AR}

\subsubsection{A funkcionális programozás előnyei}

\subsubsection{A funkcionális programozás nehézségei}

\subsection{PureScript}

\subsection{Vue.js}

%******************példa képek:

%\begin{figure}
%\centering
%\includegraphics[width=0.7\textwidth]{org}
%\caption{Organogram.}
%\end{figure}

%\begin{wrapfigure}{r}{0.4\textwidth}
%\includegraphics[width=0.4\textwidth]{test1}
%\caption{Értékteremtő lánc diagram.}
%\end{wrapfigure}

%\begin{figure}[h]
%\centering
%\begin{subfigure}{0.45\textwidth}
%\includegraphics[width=8cm]{pelda1} 
%\end{subfigure}
%\begin{subfigure}{0.4\textwidth}
%\includegraphics[width=8cm]{pelda2}
%\end{subfigure}
%\caption{A háló $t_1$ tüzelése előtt és után}
%\end{figure}

%******************felsorolas
%\begin{definíció}
%A Petri háló ($PN$) egy ($P, T, E, w, m_0$) rendszer, ahol
%\begin{enumerate}
%\item $P$ a helyek véges halmaza,
%\item $T$ a tranzíciók véges halmaza feltéve, hogy $P \cap T = \emptyset$,
%\item $E \subseteq (P \times T) \cup (T \times P)$ az élek véges halmaza,
%\item $ w: E\rightarrow N^+$ a súlyfüggvény,
%\item $ m: P\rightarrow N$ a token-eloszlás függvény,
%\item $ m_0$ a kezdeti token eloszlás.
%\end{enumerate}
%A Petri háló struktúrát $(P, T, E, w)$ módon jelöljük. (\textit{Peterson., 1981})
%\end{definíció}

\newpage
\section{Eredmények és azok értékelése}

\newpage
\section{Továbbfejlesztési lehetőségek}

A vizsgált és specifikációk készítésekor használt modellezési nyelvek (ARIS, BPMN, UML, Petri-hálók) közül én az ARIS-t találtam jelen feladathoz a legmegfelelőbbnek. A különböző diagramtípusok integrációja és azok rendszerezett kezelése mindenképp kiemeli a többi közül. Ezen felül leginkább a folyamatok leírásakor az általa kínált információtöbbletet is hasznosnak találtam, véleményem szerint a többinél nagyobb kifejezőerővel rendelkezik ez az eszköz. 

\newpage
\section*{Összefoglalás}
\addcontentsline{toc}{section}{Összefoglalás}

A mai korszerű vállalatoknak versenyképességük megőrzése céljából elengedhetetlenné vált üzleti szoftverek alkalmazása. Dolgozatom célkitűzése egy debreceni kereskedelmi profilú kisvállalkozás folyamatainak, működésének megismerése, modellezése és a számára készülő vállalatirányítási rendszer üzleti követelmény specifikációjának elkészítése volt. 


\addcontentsline{toc}{section}{Irodalomjegyzék}
%\setlength{\parindent}{0mm}

%\nocite{*}
%\bibliographystyle{agsm}

\newpage
\renewcommand{\refname}{Irodalomjegyzék}
\begin{thebibliography}{99}

\bibitem{CS}
C. Scalfani. Functional Programming Made Easier: A Step-by-Step Guide. 2021.

\bibitem{PE}
P. Freeman. PureScript by Example. 2014 -2017.
\url{https://book.purescript.org/}

\end{thebibliography}

\end{document}
